% Generated. Do not edit.
% Projection scriptural-date-location/evidence, version 1: the Scriptural Date and
% Location page, generated whole from research/chronology.toml, the Scripture
% chronology corpus and the source library by tools/tpt document-types refresh.
% document: liturgy/roman-rite/1962/propers/temporal/47-seventh-after-pentecost
\section*{Appendix: Evidence for the Scriptural Dates}
\label{sec:date-location-evidence}
{\small For each passage, as the Scriptural Date and Location table names it: what is dated, how the corpus weighs each answer, the source's own words, and the sources cited. A subject the corpus does not date is reported as such and given no date.\par}
{\small
\par\smallskip\noindent\textit{Communion: Ps. 30:3 (Vg.; Heb. 31)}\par\nobreak
\begin{itemize}[leftmargin=1.2em,itemsep=1pt,topsep=1pt,parsep=0pt]
\item \textbf{Historical setting}. No date: the corpus holds a date for this that it reports from another edition by exception, which is not printed until the corpus records the source's own qualification of it (research pending).
\item \textbf{Historical setting} (preferred). David's flight from the face of Absalom: When he fled from the face of his son Absalom (The Holy Bible, Douay-Rheims (Challoner revision), Project Gutenberg eBook 1581: Ps. 3:1).
\item \textbf{Composition} (preferred). Psalms: Before c.~165 B.C. (``before the Maccabean period, around 165 B.C.; no individual psalm can be dated securely''; Catholic critical chronology; New American Bible, Revised Edition --- USCCB web state (2026), Psalms introduction).
\end{itemize}
\par\smallskip\noindent\textit{Gradual: Ps. 33:12 (Vg.; Heb. 34), Ps. 33:6 (Vg.; Heb. 34)}\par\nobreak
\begin{itemize}[leftmargin=1.2em,itemsep=1pt,topsep=1pt,parsep=0pt]
\item \textbf{Traditional attribution} (disputed). David's reign: B.C.~1055--1015 (``reigned from 1055 to 1015 B.C.''; The Catholic Encyclopedia, Volume 4 (New York, 1908), article 'David, King,' printed p. 642, opening paragraph). Place: Hebron, then Jerusalem (The Holy Bible, Douay-Rheims (Challoner revision), Project Gutenberg eBook 1581: 2 Kings 5:5).
\item \textbf{Superscription setting} (preferred). David's flight to Achis, king of Geth: A.M.~2944 (Haydock edition reporting Ussher) (``A. M. 2944, A. C. 1060''; Douay Rheims Edition of the Holy Catholic Bible with a Comprehensive Catholic Commentary, William Peter Feeney 2014 Memorial Edition, 1K.21.1-15; Douay Rheims Edition of the Holy Catholic Bible with a Comprehensive Catholic Commentary, William Peter Feeney 2014 Memorial Edition, Ps.70.1). Place: Geth, at the Philistine court (The Holy Bible, Douay-Rheims (Challoner revision), Project Gutenberg eBook 1581: Ps. 33:1, 1 Kings 21:10, 21:15).
\item \textbf{Composition}. As under Communion: Ps. 30:3 (Vg.; Heb. 31).
\end{itemize}
\par\smallskip\noindent\textit{Alleluia: Ps. 46:2 (Vg.; Heb. 47)}\par\nobreak
\begin{itemize}[leftmargin=1.2em,itemsep=1pt,topsep=1pt,parsep=0pt]
\item \textbf{Historical setting}. No date: the chronology corpus has not yet dated this (research pending).
\item \textbf{Composition} (preferred). The Psalms of the Sons of Korah: Between the days of Isaias and the return from exile (The Catholic Encyclopedia, Volume 12 (1911)).
\end{itemize}
\par\smallskip\noindent\textit{Introit: Ps. 46:2 (Vg.; Heb. 47), Ps. 46:3 (Vg.; Heb. 47)}\par\nobreak
\begin{itemize}[leftmargin=1.2em,itemsep=1pt,topsep=1pt,parsep=0pt]
\item \textbf{Historical setting}. As under Alleluia: Ps. 46:2 (Vg.; Heb. 47).
\item \textbf{Composition}. As under Alleluia: Ps. 46:2 (Vg.; Heb. 47).
\end{itemize}
\par\smallskip\noindent\textit{Offertory: Dan. 3:40}\par\nobreak
\begin{itemize}[leftmargin=1.2em,itemsep=1pt,topsep=1pt,parsep=0pt]
\item \textbf{Historical setting} (preferred). The seventy years of servitude to Babylon: ``seventy years'' (duration) (The Holy Bible, Douay-Rheims (Challoner revision), Project Gutenberg eBook 1581: Jer. 25:11, 25:12). Place: Judah and Babylon (The Holy Bible, Douay-Rheims (Challoner revision), Project Gutenberg eBook 1581: Jer. 25:11, 25:12).
\item \textbf{Composition} (preferred). Prayer of Azarias and Song of the Three Children: Most likely long after the Exile (The Catholic Encyclopedia, Volume 4 (New York, 1908)).
\end{itemize}
\par\smallskip\noindent\textit{Gospel: Matt. 7:15--21}\par\nobreak
\begin{itemize}[leftmargin=1.2em,itemsep=1pt,topsep=1pt,parsep=0pt]
\item \textbf{Narrated event}. No date: the chronology corpus has not yet dated this (research pending).
\item \textbf{Composition} (disputed). The Gospel of St Matthew: c.~A.D.~38--45 (``about A.D. 38-45''; The Catholic Encyclopedia, Volume 10 (1911); The Catholic Encyclopedia, Volume 14 (New York, 1912)); c.~A.D.~40--42 (``about the year 40-42''; The Catholic Encyclopedia, Volume 10 (1911)); A.D.~40--45 (``the years 40-45''; The Catholic Encyclopedia, Volume 10 (1911)); c.~A.D.~60--68 (``about the year 60-68''; The Catholic Encyclopedia, Volume 10 (1911)); c.~A.D.~64--67 (``about the years 64-67''; The Catholic Encyclopedia, Volume 10 (1911)); c.~A.D.~50 (``about the year 50''; The Catholic Encyclopedia, Volume 14 (New York, 1912)). Place: Semitic original associated with Palestine (The Catholic Encyclopedia, Volume 10 (1911)).
\end{itemize}
\par\smallskip\noindent\textit{Epistle: Rom. 6:19--23}\par\nobreak
\begin{itemize}[leftmargin=1.2em,itemsep=1pt,topsep=1pt,parsep=0pt]
\item \textbf{Composition} (disputed). The Epistle to the Romans: A.D.~57 (``(Romans), 57''; The Catholic Encyclopedia, Volume 11 (New York, 1911)); A.D.~58 (``the Epistle to the Romans (58)''; The Catholic Encyclopedia, Volume 5 (New York, 1909)).
\end{itemize}
\par}
